% !TEX program = lualatex
\documentclass[12pt, a4paper]{article}


\usepackage{fontspec}
\setmainfont{Bookman Old Style}%{CMU Serif}

\usepackage{csquotes}
\usepackage[english, russian, ukrainian]{babel}

\usepackage{amsmath, amssymb}
\usepackage{geometry}
\geometry{margin=2.5cm}
\usepackage{hyperref}
\hypersetup{colorlinks=true, linkcolor=blue, urlcolor=blue, citecolor=green!50!black}
\usepackage{microtype}
%\usepackage{parskip}
\usepackage{booktabs}
\usepackage{enumitem}
\usepackage{xcolor}
\usepackage{mdframed}
\usepackage{tabularray}
\UseTblrLibrary{booktabs}
\definecolor{boxbg}{RGB}{245,245,250}
\definecolor{accent}{RGB}{30,80,160}
\definecolor{planckblue}{RGB}{31,119,180}
\definecolor{theoryred}{RGB}{214,39,40}

\usepackage{pgfplots}
\pgfplotsset{compat=1.18}
\usepackage{pgfplotstable}

\usepackage{indentfirst}

\usepackage[backend=biber,
            style=phys,         % або numeric / authoryear залежно від вимог
            sorting=none,       % сортування в порядку цитування
            maxnames=4,         % якщо авторів більше 4, скорочувати до et al.
            eprint=true,
            doi=true]{biblatex}

\addbibresource{references.bib}

\newmdenv[backgroundcolor=boxbg, linecolor=accent, linewidth=1pt,
          innerleftmargin=10pt, innerrightmargin=10pt,
          innertopmargin=8pt, innerbottommargin=8pt]{keybox}


\title{\bfseries Конспект: від інфлатона до CMB}
\author{}
\date{}

% ---------------------------------------------------------------
% Planck 2018 binned TT data
% Columns: ell  Dl[muK^2]  err_minus  err_plus
% Source: Planck Collaboration 2018, arXiv:1807.06209
% ---------------------------------------------------------------
\pgfplotstableread{plankdata.txt}\planckdata

% ---------------------------------------------------------------
% LCDM theory curve (CAMB, Planck 2018 best-fit parameters)
% H0=67.36, Ob h^2=0.02237, Oc h^2=0.1200,
% tau=0.0544, As=2.1e-9, ns=0.9649
% ---------------------------------------------------------------
\pgfplotstableread{LCDM.txt}\theorydata

\begin{document}
\maketitle
\tableofcontents
\newpage

%% --------------------------------------------------------
%% Пролог
%% --------------------------------------------------------

Цей конспект будує міст між двома точками: \emph{інфлатоном}~---
гіпотетичним скалярним полем, яке керувало першими
${\sim}10^{-32}$~с існування Всесвіту~--- і \emph{реліктовим випромінюванням} (CMB),
що є найдетальнішою спостережуваною картиною того, що відбулося.
Між цими точками лежить ланцюжок фізики, який ми будуємо
послідовно, крок за кроком.

\textbf{Частина~I. Фон.} Розділи~\ref{sec:friedmann}–\ref{sec:inflaton}
закладають геометричну та динамічну основу: рівняння Фрідмана описують
однорідний і ізотропний Всесвіт, що розширюється; інфлатон і його
потенціал $V(\phi)$ пояснюють, чому це розширення на початку було
прискореним (інфляція). Цей ідеалізований, повністю симетричний фон є
\emph{нульовим наближенням}, на якому тримається вся подальша фізика.

\textbf{Частина~II. Збурення.} Розділи~\ref{sec:metric}-\ref{sec:quantum}
переходять від симетричного фону до реального світу. Квантові флуктуації
інфлатона~--- неминучі за принципом невизначеності~--- стають
\emph{насінням} усієї великомасштабної структури. Тензорні збурення
метрики народжують первинні гравітаційні хвилі. Розділ~\ref{sec:metric}
формалізує мову, якою описуються ці збурення: SVT-розклад і
конформно-ньютонівське калібрування.

\textbf{Частина~III. Спостереження.} Розділи~\ref{sec:reheating}–\ref{sec:pol}
стежать за долею збурень після інфляції: реогрів перетворює поле на
гарячу плазму; акустичні осциляції цієї плазми залишають відбиток у
вигляді піків на спектрі CMB і масштабу BAO; поляризаційні E- і B-моди
несуть інформацію про скалярні та тензорні збурення відповідно.

\textbf{Частина~IV. Задача.} Розділи~\ref{sec:chaotic}–\ref{sec:broken}
виходять за межі стандартної моделі. Телескоп JWST виявив масивні
галактики при $z \sim 10$–$12$, існування яких важко узгодити зі
стандартним степенево-законним первинним спектром. Ми розглядаємо
механізм локального \emph{зламу спектра} через особливість у потенціалі
інфлатона і формулюємо програму перевірки цієї ідеї за допомогою
CAMB/CLASS та байєсівського аналізу (MCMC).

\bigskip
\noindent
\textit{Передумови.} Конспект розрахований на читача, знайомого із
загальною теорією відносності на рівні тензора Рімана і рівнянь
Айнштайна, та з квантовою теорією поля на рівні канонічного
квантування. Знання космології не передбачається.

\newpage

%% --------------------------------------------------------
\section{Від рівнянь Айнштайна до рівнянь Фрідмана}
\label{sec:friedmann}
%% --------------------------------------------------------

Космологія описує Всесвіт як єдиний фізичний об'єкт. Її відправний
пункт~--- принцип космологічної однорідності та ізотропії: на
достатньо великих масштабах ($\gtrsim 100\,\mathrm{Mpc}$) розподіл
матерії та геометрія простору-часу не виділяють жодного напрямку і
жодної точки. Цей принцип фіксує єдиний клас метрик~---
Фрідмана--Леметра--Робертсона--Вокера (FLRW)~--- і зводить
рівняння Айнштайна до двох звичайних диференціальних рівнянь для
одного невідомого: масштабного фактора $a(t)$.

Метрика FLRW у \emph{космологічному часі}~$t$ для плоского Всесвіту
($k = 0$):
\begin{equation}\label{eq:FLRW}
  ds^2 = -dt^2 + a(t)^2\,\delta_{ij}\,dx^i dx^j.
\end{equation}
Підставляючи~\eqref{eq:FLRW} у рівняння Айнштайна
$G_{\mu\nu} + \Lambda g_{\mu\nu} = 8\pi G\,T_{\mu\nu}$
для ідеального плину
$T^{\mu\nu} = \mathrm{diag}(\rho,\,p,\,p,\,p)$
і обчислюючи ненульові компоненти тензора Айнштайна
\begin{equation*}
  G^0{}_0 = -3H^2,
  \qquad
  G^i{}_j = -(2\dot H + 3H^2)\,\delta^i{}_j,
  \qquad
  H \equiv \frac{\dot a}{a},
\end{equation*}
отримуємо \textbf{рівняння Фрідмана}:

\begin{keybox}
\textbf{Перше рівняння Фрідмана} ($(00)$-компонента):
\begin{equation}\label{eq:F1}
  H^2 = \left(\frac{\dot a}{a}\right)^2
      = \frac{8\pi G}{3}\,\rho + \frac{\Lambda}{3}
\end{equation}
\textbf{Друге рівняння Фрідмана} ($(ii)$-компонента, «прискорення»):
\begin{equation}\label{eq:F2}
  \frac{\ddot a}{a} = -\frac{4\pi G}{3}\,(\rho + 3p) + \frac{\Lambda}{3}
\end{equation}
\textbf{Рівняння збереження} ($\nabla_\mu T^{\mu\nu} = 0$):
\begin{equation}\label{eq:conservation_low}
  \dot\rho + 3H(\rho + p) = 0
\end{equation}
\end{keybox}

\noindent
\textit{Зауваження:}
\begin{itemize}
  \item Рівняння збереження не є незалежним: воно автоматично
        слідує з двох рівнянь Фрідмана через тотожність Б'янкі
        $\nabla_\mu G^{\mu\nu} = 0$.
  \item Космологічна стала $\Lambda$ еквівалентна ідеальній рідині
        з $\rho_\Lambda = \Lambda/(8\pi G)$, $p_\Lambda = -\rho_\Lambda$,
        тобто $w = -1$.
  \item Відхилення від плоскості ($k = \pm 1$) додає член
        $-k/a^2$ у перше рівняння; Planck обмежує
        $|\Omega_k| < 0.005$, тому $k = 0$~--- відмінне наближення.
\end{itemize}

%% --------------------------------------------------------
\subsection*{Рівняння стану і закони розширення}
%% --------------------------------------------------------

Кожна компонента Всесвіту характеризується \textbf{рівнянням стану}
$w \equiv p/\rho$. Рівняння збереження~\eqref{eq:conservation_low}
при сталому $w$ інтегрується точно:
\[
  \dot\rho + 3\frac{\dot a}{a}(1+w)\rho = 0
  \;\Rightarrow\;
  \rho \propto a^{-3(1+w)}.
\]
Підставляємо у перше рівняння Фрідмана $H^2 \propto \rho$:
\[
  \left(\frac{\dot a}{a}\right)^2 \propto a^{-3(1+w)}
  \;\Rightarrow\;
  a(t) \propto t^{\,\frac{2}{3(1+w)}}.
\]
Для випромінювання $\rho \propto a^{-4}$: додатковий множник $a^{-1}$
порівняно з матерією~--- це червоне зміщення фотонів ($E_\gamma \propto 1/a$).

\begin{center}
\begin{tblr}{
  colspec={lllll},
  row{1} = {c, m, font=\bfseries}
}
\toprule
Епоха & Джерело & $w$ & $\rho(a)$ & $a(t)$ \\
\midrule
Інфляція    & Інфлатон              & $\approx -1$ & $\approx\mathrm{const}$ & $\propto e^{Ht}$ \\
Радіаційна  & Фотони, нейтрино      & $1/3$        & $\propto a^{-4}$        & $\propto t^{1/2}$ \\
Матеріальна & Темна матерія+баріони & $0$          & $\propto a^{-3}$        & $\propto t^{2/3}$ \\
Сьогодні    & Темна енергія         & $\approx -1$ & $\approx\mathrm{const}$ & $\propto e^{Ht}$ \\
\bottomrule
\end{tblr}
\end{center}

%% --------------------------------------------------------
\subsection*{Червоний зсув}
%% --------------------------------------------------------

Спостережна кінематика космосу пов'язана з $a(t)$ через
\textbf{космологічний червоний зсув}~$z$. Фотон, випромінений
у момент $t_e$, до спостерігача ($t_0$) приходить розтягненим:
\begin{equation}\label{eq:redshift}
  1 + z \equiv \frac{\lambda_{\rm obs}}{\lambda_{\rm em}}
             = \frac{a(t_0)}{a(t_e)}.
\end{equation}
Нормуємо $a(t_0) \equiv 1$, тоді $a = 1/(1+z)$.
Характерні значення:
\begin{align*}
  z &= 0\ \text{(сьогодні)}, \\
  z &\approx 0.3\ \text{(прискорення розширення)}, \\
  z &\approx 1100\ \text{(рекомбінація)}, \\
  z &\sim 10^{10}\ \text{(нуклеосинтез)}.
\end{align*}

%% --------------------------------------------------------
\subsection*{Параметри густини і модель $\Lambda$CDM}
%% --------------------------------------------------------

Зручно нормувати густину кожної компоненти на
\textbf{критичну густину}~--- значення, за якого Всесвіт плаский:
\begin{equation}\label{eq:rho_crit}
  \rho_{\rm cr}(t) \equiv \frac{3H^2(t)}{8\pi G},
  \qquad
  \Omega_i \equiv \frac{\rho_i}{\rho_{\rm cr}}.
\end{equation}
Тоді перше рівняння Фрідмана~\eqref{eq:F1} переписується як:
\[
  \sum_i \Omega_i = 1 \quad (k=0),
\]
а параметр Хаббла через $z$:
\begin{equation}\label{eq:Hz}
  H(z) = H_0\sqrt{\Omega_r(1+z)^4 + \Omega_m(1+z)^3 + \Omega_\Lambda},
\end{equation}
де кожен член~--- внесок своєї компоненти з таблиці вище.

Стандартна космологічна модель $\Lambda$CDM описує Всесвіт
чотирма компонентами:

\begin{center}
\begin{tblr}{colspec={llll}, row{1}={font=\bfseries}}
\toprule
Компонента & $\Omega_i$ & $w$ & Частка сьогодні \\
\midrule
Баріонна матерія          & $\Omega_b$       & $0$    & ${\approx}5\%$  \\
Холодна темна матерія     & $\Omega_c$       & $0$    & ${\approx}27\%$ \\
Темна енергія ($\Lambda$) & $\Omega_\Lambda$ & $-1$   & ${\approx}68\%$ \\
Випромінювання            & $\Omega_r$       & $1/3$  & ${\approx}0\%$  \\
\bottomrule
\end{tblr}
\end{center}

З погляду спостережної космології, модель повністю параметризується
мінімальним набором із шести чисел\footnote{Фізичні густини баріонів
$\Omega_b h^2$ та темної матерії $\Omega_c h^2$, кутовий масштаб
звукового горизонту $\theta_\ast$, оптична глибина реіонізації $\tau$,
амплітуда $A_s$ і спектральний індекс $n_s$ первинних збурень.}.
Їхні найдостовірніші значення зафіксовані місією Planck~2018
\cite{planck_2018}.

Попри концептуальну простоту, $\Lambda$CDM блискуче описує
кутовий спектр CMB, великомасштабну структуру (LSS),
баріонні акустичні осциляції (BAO) та діаграму Хаббла для наднових~Ia.
Однак вона є теорією \emph{еволюції}, а не \emph{початкових умов}:
параметри $A_s$, $n_s$ приймаються феноменологічно.
Їхнє фізичне обґрунтування дає теорія інфляції~--- предмет наступного розділу.

%% --------------------------------------------------------
\subsection*{Конформний час і горизонт Хаббла}
%% --------------------------------------------------------

Для опису еволюції збурень зручніше перейти від
космологічного часу $t$ до \emph{конформного часу}
\begin{equation*}
  d\eta = \frac{dt}{a},
\end{equation*}
у якому метрика FLRW набуває конформно-плаского вигляду:
\begin{equation*}
  ds^2 = a(\eta)^2\bigl(-d\eta^2 + \delta_{ij}\,dx^i dx^j\bigr).
\end{equation*}
У цих координатах розширення Всесвіту повністю «заховане»
у множнику $a(\eta)$~--- координати окремих галактик із часом
не змінюються\footnote{Це справедливо лише для галактик,
що беруть участь у загальному космологічному розширенні
(хаббловському потоці). Галактики з власним рухом відносно
нього~--- наприклад, Андромеда (M31), що наближається до
Чумацького Шляху зі швидкістю ${\sim}110\,\mathrm{km/s}$~---
мають координати, що змінюються з часом. На космологічних
масштабах такі \emph{пекулярні швидкості} малі порівняно
з темпом розширення і зазвичай ними нехтують.}.

\begin{keybox}
Перевага конформного часу полягає в тому, що хвильове
рівняння для будь-якої моди збурень у цих змінних точно
відповідає рівнянню гармонічного осцилятора у пласкому
просторі Мінковського~--- множник $a(\eta)$ виноситься,
і розширення Всесвіту проявляється лише через ефективний
часозалежний потенціал. Це робить конформний час природною
мовою теорії збурень, яка розвивається у розділах
\ref{sec:metric}–\ref{sec:quantum}.
\end{keybox}

Позначимо похідну за $\eta$ штрихом і введемо конформний
параметр Хаббла $\mathcal{H} \equiv a'/a = aH$. Рівняння
Фрідмана у конформних змінних:
\begin{equation*}
  \mathcal{H}^2 = \frac{8\pi G}{3}\,a^2\bar{\rho}
                + \frac{\Lambda}{3}\,a^2,
  \qquad
  \mathcal{H}' = -\frac{4\pi G}{3}\,a^2(\bar{\rho} + 3\bar{p})
               + \frac{\Lambda}{3}\,a^2,
\end{equation*}
де $\bar{\rho}(\eta)$ та $\bar{p}(\eta)$~--- однорідні фонові
густина та тиск. Комбінуючи ці два рівняння, отримуємо чисто
кінематичне співвідношення, що зв'язує темп зміни горизонту
з рівнянням стану:
\begin{equation}\label{eq:calH_prime}
  \mathcal{H}' = \mathcal{H}^2\!\left[1 - \tfrac{3}{2}(1+w)\right],
  \qquad w \equiv \frac{\bar{p}}{\bar{\rho}}.
\end{equation}
Знак правої частини~\eqref{eq:calH_prime} визначає динаміку
горизонту: при $w > -1/3$ (стандартне розширення) фізичний
радіус Хаббла $R_H = c/H$ зростає, а при $w < -1/3$
(інфляція)~--- стискається. Фізично $R_H$~--- це масштаб,
на якому швидкість розширення Всесвіту дорівнює швидкості
світла; у теорії збурень він відокремлює режими, що
коливаються ($k \gg \mathcal{H}$), від «заморожених»
($k \ll \mathcal{H}$)\footnote{Це поняття динамічного
горизонту не слід плутати з горизонтом частинок, який
визначає абсолютну причинну зв'язаність від початку часів.}.

Розв'язуючи~\eqref{eq:calH_prime} для сталого $w$, отримуємо
закон розширення $a(\eta)$ для кожної космологічної епохи:

\begin{center}
\begin{tblr}{
  colspec={lllll},
  row{1} = {c, m, font=\bfseries}
}
\toprule
Епоха & Джерело & $w$ & $a(\eta)$ & $a(t)$ \\
\midrule
Інфляція    & Інфлатон
            & $\approx -1$
            & $\propto -\dfrac{1}{H\eta}\ \ (\eta < 0)$\footnotemark
            & $\propto e^{Ht}$ \\
Радіаційна  & Фотони, нейтрино
            & $1/3$
            & $\propto \eta$
            & $\propto t^{1/2}$ \\
Матеріальна & {Темна матерія \\ + баріони}
            & $0$
            & $\propto \eta^2$
            & $\propto t^{2/3}$ \\
Сьогодні    & Темна енергія
            & $\approx -1$
            & $\propto -\dfrac{1}{H\eta}\ \ \text{(асимптотика)}$
            & $\propto e^{Ht}$ \\
\bottomrule
\end{tblr}
\end{center}
\footnotetext{Під час інфляції $a(t)=a_0 e^{Ht}$, тому
$\eta(t)=\int_{-\infty}^{t}\frac{dt'}{a_0 e^{Ht'}}
=-\frac{1}{a(t)H}=-\frac{1}{\mathcal{H}}$.
Оскільки $a(t)H>0$ завжди, $\eta<0$ протягом усієї інфляції;
зручно покласти $\eta\to 0^-$ на її кінці, тобто вся епоха~---
$\eta\in(-\infty,0)$.
Конформний горизонт Хаббла $\mathcal{H}^{-1}=-\eta=1/(aH)$
при цьому \emph{спадає}: $a$ росте при $H\approx\mathrm{const}$,
що є кінематичним визначенням інфляції.}

%% --- Місток до §2 ---
$\Lambda$CDM є теорією \emph{еволюції}, а не
\emph{початкових умов}: вона не виводить з перших принципів (\emph{ab initio})
ні надзвичайної просторової пласкості Всесвіту, ні природи
і спектру первинних флуктуацій густини~--- ці параметри
($A_s$, $n_s$, про них йтиметься далі) приймаються як феноменологічні вхідні дані,
що визначаються з спостережень. Їхнє фізичне обґрунтування
забезпечує теорія космічної інфляції~--- предмет наступного розділу.

%% --------------------------------------------------------
\section{Інфлатон і його потенціал}
\label{sec:inflaton}
%% --------------------------------------------------------

Як було зазначено, стандартна космологія стикається
з двома фундаментальними проблемами. \emph{Проблема горизонту:}
протилежні ділянки неба на картах CMB мають однакову температуру
з точністю до $10^{-5}$, хоча в момент рекомбінації вони не могли
перебувати у причинному контакті. \emph{Проблема пласкості:}
просторова кривина Всесвіту сьогодні надзвичайно мала
($|\Omega_k| < 0.005$), що вимагає надточного тонкого підбору
початкових умов. Обидві проблеми розв'язуються одним механізмом:
якщо у ранньому Всесвіті існував короткий період прискореного
розширення (\emph{інфляція}), то сьогоднішній спостережуваний
Всесвіт є крихітною та сильно вирівняною областю набагато більшого
початкового об'єму.

З погляду релятивістської гідродинаміки, для запуску такого прискореного
розширення ($\ddot{a} > 0$) необхідне середовище з екзотичним рівнянням
стану, що задовольняє умову $p < -\rho c^2 / 3$, тобто має сильний
від'ємний тиск. Звичайна речовина або випромінювання на це не здатні.

Найпростішою реалізацією такого механізму в теорії поля є однопольова
модель, де прискорене розширення забезпечується однорідним скалярним
полем \emph{інфлатоном} $\phi(t)$, яке повільно «котиться» вздовж
пологого потенціалу $V(\phi)$. Фізичну аналогію можна провести з
рухом важкої кульки по пологому схилу за наявності сильного
в'язкого тертя: кулька рухається повільно і рівномірно, не
набираючи кінетичної енергії. Завдяки просторовій однорідності
фону, на цьому етапі просторовими градієнтами поля можна знехтувати.

Дія для скалярного поля $\phi$ у просторі FLRW:
\begin{equation}\label{eq:inflaton_action}
  S = \int d^4x\,\sqrt{-g}\left[\frac{1}{2}g^{\mu\nu}\partial_\mu\phi\,\partial_\nu\phi - V(\phi)\right].
\end{equation}
Для однорідного поля $\phi(t)$ маємо $\sqrt{-g} = a^3$ і $g^{00} = -1$,
тому підінтегральний вираз зводиться до:
\begin{equation*}
  \mathcal{L} = a^3\!\left[\frac{\dot{\phi}^2}{2} - V(\phi)\right].
\end{equation*}
Рівняння Ейлера--Лагранжа $\frac{d}{dt}\frac{\partial\mathcal{L}}{\partial\dot{\phi}} - \frac{\partial\mathcal{L}}{\partial\phi} = 0$ дають:
\begin{equation*}
  \frac{d}{dt}\!\left(a^3\dot{\phi}\right) + a^3 V'(\phi) = 0,
\end{equation*}
що після ділення на $a^3$ зводиться до:
\begin{equation}\label{eq:klein_gordon}
    \ddot{\phi} +3H\dot{\phi} + V'(\phi) = 0,
\end{equation}
де $3H\dot{\phi}$ виникає з диференціювання $a^3$:
$\frac{d}{dt}(a^3\dot{\phi}) = a^3(\ddot{\phi} + 3H\dot{\phi})$~---
це і є те саме «хабблівське тертя», яке фізично відіграє роль
в'язкого опору у аналогії з кулькою на схилі.

Варіюючи ту саму дію~\eqref{eq:inflaton_action} по метриці,
отримуємо тензор енергії-імпульсу скалярного поля; для
однорідного $\phi(t)$ він набуває вигляду ідеальної рідини
з густиною і тиском:
\begin{equation}\label{eq:inflaton_rho_p}
  \bar{\rho}_\phi = \frac{\dot{\phi}^2}{2} + V(\phi),
  \qquad
  \bar{p}_\phi   = \frac{\dot{\phi}^2}{2} - V(\phi).
\end{equation}


\noindent
\textit{Зауваження про нотацію.} У цьому та наступних розділах
штрих у виразах типу $V'(\phi)$ завжди означає похідну за
скалярним полем $\phi$, тобто $V' \equiv dV/d\phi$. Це відрізняється
від §\,\ref{sec:friedmann}, де штрих позначав похідну за конформним
часом~$\eta$.

З виразів~\eqref{eq:inflaton_rho_p} безпосередньо випливає умова
інфляції. Розглянемо перше рівняння Фрідмана~\eqref{eq:F1} для
інфлатона як єдиного джерела енергії, де зручно ввести
\emph{редуковану масу Планка} $M_\mathrm{Pl} \equiv (8\pi G)^{-1/2}$:
\begin{equation}\label{eq:friedmann_inflaton}
  H^2 = \frac{1}{3M_\mathrm{Pl}^2}
        \left(\frac{\dot\phi^2}{2} + V(\phi)\right).
\end{equation}
Позначимо $\varepsilon \equiv \dot\phi^2/(2H^2 M_\mathrm{Pl}^2)$
(формально визначається у §\,\ref{subsec:slowroll}).
Підставляючи $H^2 M_\mathrm{Pl}^2 = (\dot\phi^2/2 + V)/3$
з~\eqref{eq:friedmann_inflaton}:
\begin{equation*}
  \varepsilon
    = \frac{\dot\phi^2/2}{2H^2 M_\mathrm{Pl}^2}
    = \frac{\dot\phi^2/2}{\tfrac{2}{3}(\dot\phi^2/2 + V)}
    = \frac{\dot\phi^2/2}{\dot\phi^2/2 + V(\phi)}.
\end{equation*}
Умова $\varepsilon \ll 1$ тому \emph{еквівалентна} умові
$\dot\phi^2 \ll V(\phi)$, і у цьому режимі:
\begin{equation*}
  H^2 \approx \frac{V(\phi)}{3M_\mathrm{Pl}^2}, \qquad
  w_\phi = \frac{\bar{p}_\phi}{\bar{\rho}_\phi}
         = \frac{\dot\phi^2/2 - V}{\dot\phi^2/2 + V} \approx -1,
\end{equation*}
тобто густина енергії залишається практично сталою
($\bar{\rho}_\phi \approx V(\phi) \approx \mathrm{const}$).
Підставляючи у перше рівняння Фрідмана, отримуємо
сталий параметр Хаббла $H \approx \mathrm{const}$ та
експоненційне розширення:
\begin{equation}\label{eq:a_exponential}
  a(t) = a_0\,e^{Ht}.
\end{equation}
Саме цей режим вирішує проблеми горизонту і пласкості: за $N \sim
50$--$60$ $e$-folds ($a$ зростає в $e^{50}$--$e^{60}$ разів)
будь-яка початкова кривина розгладжується, а причинно пов'язана
область стає набагато більшою за спостережуваний Всесвіт.

%% --------------------------------------------------------
\subsection{Параметри повільного кочення (slow-roll)}
\label{subsec:slowroll}
%% --------------------------------------------------------

Режим повільного кочення формалізується двома безрозмірними
параметрами, що вимірюють, наскільки далеко система від чистого
де-Сіттерівського розширення:
%\textbf{Параметри slow-roll через динаміку Хаббла:}
\begin{equation}\label{eq:slow_roll_def}
  \varepsilon \equiv -\frac{\dot H}{H^2}
             = \frac{\dot\phi^2}{2H^2 M_\mathrm{Pl}^2},
  \qquad
  \eta_{\mathrm{H}} \equiv -\frac{\ddot{\phi}}{H\dot{\phi}},
\end{equation}
де $M_\mathrm{Pl}$~--- редукована маса Планка (означена вище).
Параметр $\varepsilon$ вимірює відносну зміну $H$ за один $e$-fold;
$\eta_{\mathrm{H}}$~--- відносне прискорення поля.


З означення~\eqref{eq:slow_roll_def} маємо
$|\eta_\mathrm{H}| = |\ddot\phi|/(H|\dot\phi|)$, тому умова
$|\eta_\mathrm{H}| \ll 1$ еквівалентна $|\ddot\phi| \ll H|\dot\phi|$,
що є достатнім для виконання $|\ddot\phi| \ll 3H|\dot\phi|$.
У цьому наближенні рівняння~\eqref{eq:klein_gordon} зводиться
до першого порядку:
\begin{equation}\label{eq:sr_attractor}
  3H\dot{\phi} \approx -V'(\phi).
\end{equation}
Підставляючи~\eqref{eq:sr_attractor} у $\varepsilon$, виражаємо
параметри безпосередньо через геометрію потенціалу:
\begin{equation}\label{eq:slow_roll_pot}
  \varepsilon \approx \epsilon_V
    \equiv \frac{M_\mathrm{Pl}^2}{2}\!\left(\frac{V'}{V}\right)^2,
  \qquad
  \eta_{\mathrm{H}} \approx \eta_V
    \equiv M_\mathrm{Pl}^2\,\frac{V''}{V}.
\end{equation}
\textit{Виведення $\epsilon_V$:} з~\eqref{eq:sr_attractor}
маємо $\dot\phi \approx -V'/(3H)$, тому
\begin{equation*}
  \varepsilon
    = \frac{\dot\phi^2}{2H^2 M_\mathrm{Pl}^2}
    \approx \frac{(V')^2}{18 H^4 M_\mathrm{Pl}^2}
    \overset{H^2\approx V/(3M_\mathrm{Pl}^2)}{=\!=\!=\!=}
    \frac{M_\mathrm{Pl}^2}{2}\!\left(\frac{V'}{V}\right)^2
    \equiv \epsilon_V.
\end{equation*}
Умова інфляції~--- $\varepsilon < 1$ (еквівалентно $\ddot{a} > 0$).
Для її тривалості необхідно $\varepsilon \ll 1$ та $|\eta_V| \ll 1$:
потенціал має бути пологим і майже лінійним на потрібній ділянці.
Інфляція завершується, коли поле прискорюється настільки, що
$\varepsilon(\phi_{\mathrm{end}}) = 1$.

%% --------------------------------------------------------
\subsection{Геометрія потенціалу та обмеження спостережень}
%% --------------------------------------------------------

Спостереження CMB дають нам інформацію не про весь потенціал
$V(\phi)$, а лише про надзвичайно вузький його проміжок. Масштаби,
зафіксовані на картах анізотропії Planck, виходили за горизонт
протягом $\sim\!7$--$8$ $e$-folds із загальних $50$--$60$. Іншими
словами, кожне кутове масштабне відхилення температури CMB
відповідає певному значенню $\phi$ під час інфляції, і вся ця
«пам'ять» закодована лише у вузькій ділянці потенціалу~---
тій, що залишила слід на кутовому спектрі анізотропій
(див.\ рис.~\ref{fig:cmb_tt}).

%\begin{figure}[htbp]
%\centering
%\includegraphics[width=\linewidth]{planck_cmb_map.png}
%\caption{Карта анізотропій CMB від місії Planck~2018
%         \cite{planck_2018}. Кольори відображають відхилення
%         температури від середнього значення $T_0 = 2.725\,\mathrm{K}$
%         в діапазоні $\pm 300\,\mu\mathrm{K}$. Саме ця карта
%         є спостережуваним відбитком первинних флуктуацій
%         інфлатона.}
%\label{fig:planck_map}
%\end{figure}

Сьогодні існує понад 100 конкурентних моделей інфляції.
Більшість із них дають практично однаковий фіт для великих
кутових масштабів, де Planck найточніший. Різниця між
моделями виявляється на малих масштабах
($k \gg 1\,\mathrm{Mpc}^{-1}$), де поточні спостереження
мають значно слабші обмеження.

Таблиця~\ref{tab:inflation_models} порівнює передбачення
найпопулярніших класів моделей для двох вимірюваних
параметрів~--- $n_s$ і $r$~--- з поточними обмеженнями
Planck~2018 \cite{planck_2018} + BICEP/Keck \cite{bicep_keck_2021}.

\begin{table}[htbp]
\centering
\caption{Передбачення основних класів інфляційних моделей.
         Значення наведені для $N = 50$--$60$ $e$-folds до
         кінця інфляції. Спостережувані межі:
         $n_s = 0.9649 \pm 0.0042$,\ $r < 0.036$
         \cite{planck_2018, bicep_keck_2021}.}
\label{tab:inflation_models}
\begin{tblr}{
  colspec = {XXXXX},
  row{1}  = {font=\bfseries},
}
\toprule
Модель & Потенціал $V(\phi)$ & $n_s$ & $r$ & Статус \\
\midrule
Старобінський ($R^2$) \cite{starobinsky_1980}
  & плато & $1 - \frac{2}{N}$ & $\frac{12}{N^2}$
  & узгоджується \\
Хіггсова інфляція \cite{bezrukov_2008}
  & плато (нем.~зв'язок) & $\approx 0.965$ & $\approx 0.004$
  & узгоджується \\
Natural inflation \cite{freese_1990}
  & $1 - \cos(\phi/f)$ & $\approx 0.95$--$0.97$ & $0.01$--$0.1$
  & під питанням \\
Хаотична, $\phi^2$ \cite{linde_1983}
  & $\frac{1}{2}m^2\phi^2$ & $1 - \frac{2}{N}$ & $\frac{8}{N}$
  & виключена ($r \sim 0.13$) \\
Хаотична, $\phi^{4/3}$ \cite{linde_1983}
  & $\lambda\phi^{4/3}$ & $1 - \frac{5}{3N}$ & $\frac{16}{3N}$
  & під питанням \\
Зі сходинкою (наша модель)
  & плато + $\delta V$ & залежить від $k_{\rm break}$ & мале
  & перевіряється \\
\bottomrule
\end{tblr}
\end{table}

Якщо геометрія потенціалу на цій «непомітній» ділянці різко
змінюється~--- наприклад, з'являється локальна сходинка або ділянка
ультра-повільного кочення (USR)~--- режим slow-roll тимчасово
порушується. Флуктуації, що виходять за горизонт саме в цей момент,
зазнають різкого підсилення, формуючи \emph{зламаний спектр
потужності}. Цей механізм є центральною ідеєю нашої роботи і
детально розглядається у §\,\ref{sec:broken}. Але перш ніж перейти
до нього~--- нам потрібно зрозуміти, як саме квантові флуктуації
інфлатона перетворюються на початкові неоднорідності Всесвіту.
Саме тут інфляція є унікальною: на відміну від будь-якого
сценарію з довільними початковими умовами, вона генерує ці
неоднорідності з квантового вакууму~--- без жодного зовнішнього
задання, як неминучий наслідок принципу невизначеності.


%% --------------------------------------------------------
\section{Збурення метрики: SVT-розклад і калібрування}
\label{sec:metric}
%% --------------------------------------------------------

Рівняння Фрідмана описують ідеально симетричний фон.
Реальний Всесвіт відхиляється від нього: є флуктуації густини,
потоки речовини, гравітаційні хвилі. Щоб описати ці відхилення
кількісно~--- і зв'язати їх з CMB~--- потрібна мова малих збурень
навколо метрики FLRW. Ця мова~--- лінійна теорія збурень метрики~---
стане фундаментом для квантування збурень у наступному розділі.

Малі відхилення $\delta g_{\mu\nu}(\eta, \mathbf{x})$ від
фонового значення $\bar{g}_{\mu\nu}(\eta)$ класифікуються за
їхньою поведінкою під просторовими обертаннями. Ключовий факт:
оскільки фон є однорідним і ізотропним, збурення різних
тензорних рангів не взаємодіють між собою у лінійному порядку~---
скалярні, векторні і тензорні компоненти еволюціонують
незалежно. Це твердження відоме як SVT-розклад
\cite{stewart_walker_1974}:
\begin{equation}\label{eq:svt_decomposition}
  \delta g_{\mu\nu}
    = \delta g_{\mu\nu}^{(\mathrm{S})}
    + \delta g_{\mu\nu}^{(\mathrm{V})}
    + \delta g_{\mu\nu}^{(\mathrm{T})}.
\end{equation}
Незалежність трьох секторів означає, що кожен з них можна
вивчати окремо:
\begin{itemize}
  \item \textbf{Скалярні збурення ($\mathrm{S}$):}
        Параметризуються двома скалярними функціями
        $(\Phi, \Psi)$. Динамічно пов'язані з флуктуаціями
        густини матерії $\delta\rho$~--- це джерело
        великомасштабної структури.
  \item \textbf{Векторні збурення ($\mathrm{V}$):}
        Описують вихрові рухи і моменти імпульсу. За
        відсутності джерел вони згасають як $a^{-2}$
        \cite{mukhanov_2005} і в стандартній моделі ролі не
        відіграють. Ми вводимо їх тут для повноти
        класифікації і далі не розглядаємо.
  \item \textbf{Тензорні збурення ($\mathrm{T}$):}
        Вільні поперечні безслідові деформації просторової
        метрики~--- первинні гравітаційні хвилі. Народжуються
        під час інфляції і практично не взаємодіють з
        речовиною.
\end{itemize}

%% --------------------------------------------------------
\subsection{Конформно-ньютонівське калібрування}
%% --------------------------------------------------------

Рівняння загальної теорії відносності є загально-коваріантними: одна і та сама фізична
ситуація описується нескінченно багатьма системами координат.
Для збурень це означає наявність фіктивних (калібрувальних)
ступенів вільності. Щоб отримати однозначний фізичний зміст,
фіксуємо \emph{конформно-ньютонівське калібрування} (поздовжнє).
У ньому збурена метрика записується як:
\begin{equation}\label{eq:perturbed_metric}
  ds^2 = a(\eta)^2\Bigl[
    -(1 + 2\Phi)\,d\eta^2
    + \bigl((1 - 2\Psi)\,\delta_{ij} + h_{ij}\bigr)
      dx^i dx^j
  \Bigr],
\end{equation}
де $\Phi$ та $\Psi$~--- калібрувально-інваріантні потенціали
Бардіна \cite{bardeen_1980}, а $h_{ij}$~--- тензорне збурення.
Фізичний зміст кожного доданку:
\begin{enumerate}
  \item $\Phi$~--- ньютонівський гравітаційний потенціал:
        визначає прискорення пробних частинок і гравітаційне
        червоне зміщення фотонів CMB.
  \item $\Psi$~--- потенціал просторової кривини: описує
        локальне стиснення або розтягнення просторових
        гіперповерхонь. За відсутності анізотропного напруження
        $\Phi = \Psi$.
  \item $h_{ij}$~--- задовольняє умовам поперечності
        ($\partial^i h_{ij} = 0$) і безслідовості
        ($h^i{}_i = 0$); описує два стани поляризації
        гравітаційних хвиль $h_+$ та $h_\times$.
\end{enumerate}

%% --------------------------------------------------------
\subsection{Динаміка первинних гравітаційних хвиль}
%% --------------------------------------------------------

Тензорна компонента $h_{ij}$ щойно з'явилась у метриці~---
природно запитати, як вона еволюціонує. Підставляючи тензорну
частину~\eqref{eq:perturbed_metric} у лінеаризовані рівняння
Айнштайна без джерел анізотропного тиску, отримуємо хвильове
рівняння для кожної Фур'є-моди:
\begin{equation}\label{eq:tensor_wave_eq}
  h_{ij}'' + 2\mathcal{H}\,h_{ij}' + k^2\,h_{ij} = 0.
\end{equation}
Його структура ідентична рівнянню~\eqref{eq:fourier_field}
для скалярних флуктуацій~--- той самий осцилятор із
хабблівським тертям $2\mathcal{H}$. Тому і результат
аналогічний: моди заморожуються за горизонтом під час
інфляції, і їхній первинний спектр визначається масштабом
Хаббла \cite{starobinsky_1979}:
\begin{equation}\label{eq:tensor_power_spec}
  \mathcal{P}_h(k)
    = \frac{2}{\pi^2 M_\mathrm{Pl}^2}
      \left(\frac{H}{2\pi}\right)^{\!2}
    = \frac{H^2}{\pi^2 M_\mathrm{Pl}^2}.
\end{equation}
Множник $M_\mathrm{Pl}^{-2}$ відрізняє тензорний спектр від
скалярного~\eqref{eq:curvature_power_spec}: гравітаційні хвилі
пригнічені планківською масою, що відображає слабкість
гравітаційної взаємодії.

Відношення амплітуд тензорних і скалярних збурень~---
\emph{тензорно-скалярне відношення}~--- безпосередньо
вимірюється через B-моди поляризації CMB:
\begin{equation}\label{eq:tensor_scalar_ratio}
  r \equiv \frac{\mathcal{P}_h(k_*)}
               {\mathcal{P}_{\mathcal{R}}(k_*)}
    = 16\varepsilon,
\end{equation}
де $k_*$~--- опорний масштаб. Сучасна верхня межа
$r < 0.036$ \cite{bicep_keck_2021} обмежує енергетичний
масштаб інфляції і відсікає моделі з крутим потенціалом.
Зв'язок між $r$ і B-модами детально розглядається у
§\,\ref{sec:pol}.

Таким чином, цей розділ побудував мову: збурена метрика
параметризована потенціалами $\Phi$, $\Psi$ і тензором $h_{ij}$,
а SVT-розклад гарантує їхню незалежну еволюцію. Тепер можна
поставити центральне питання: звідки взялися самі збурення?
Відповідь~--- квантові флуктуації інфлатона~--- є темою
наступного розділу.


%% --------------------------------------------------------
\section{Квантові флуктуації і збурення}
\label{sec:quantum}
%% --------------------------------------------------------

У §\,\ref{sec:inflaton} інфлатон розглядався як класичне однорідне
поле $\bar{\phi}(t)$, а у §\,\ref{sec:metric} ми побудували мову
для опису збурень метрики. Тепер можна поставити питання про
їхнє походження. Будь-яке квантове поле неминуче має
флуктуації~--- навіть у вакуумному стані. У розширюваному просторі
де-Сіттера ці нульові коливання набувають особливої долі: унаслідок
інфляційного розширення вони «виштовхуються» за горизонт Хаббла і
там замерзають, перетворюючись на класичні просторові неоднорідності.
Саме цей механізм є джерелом початкової структури Всесвіту.

У лінійному наближенні квантові збурення поділяються на два
незалежні класи:
\begin{enumerate}
  \item \textbf{Флуктуації інфлатона} $\delta\phi(\eta, \mathbf{x})$~---
        породжують первинні скалярні збурення кривини, що стають
        насінням великомасштабної структури.
  \item \textbf{Флуктуації метрики} $h_{ij}(\eta, \mathbf{x})$~---
        тензорні збурення, що інтерпретуються як первинні
        гравітаційні хвилі.
\end{enumerate}
Два класи еволюціонують незалежно і залишають різні відбитки
у поляризації CMB (E- і B-моди відповідно~--- §\,\ref{sec:pol}).

\textit{Зауваження про квантово-класичний перехід.}
Після виходу за горизонт квантові моди «заморожуються» і
поводяться як класичні стохастичні поля~--- саме їх ми й
спостерігаємо у вигляді температурних плям CMB. Точний механізм
цього переходу (декогеренція, «стиснення» квантового стану,
вибір базису вказівників) є відкритим концептуальним питанням
і виходить за рамки цього конспекту. Проте це питання
концептуальне, а не вимірювальне: гауссова статистика заморожених
мод, яку передбачає стандартний підхід, прекрасно узгоджується зі
спостереженнями CMB~--- зокрема з відсутністю первинної
негауссовості за даними Planck~2018 \cite{planck_2018}.
Для практичних розрахунків достатньо того, що амплітуди
заморожених мод підпорядковуються гауссовій статистиці з
дисперсією~\eqref{eq:field_power_spec}.

%% --------------------------------------------------------
\subsection{Скалярні флуктуації у просторі де-Сіттера}
%% --------------------------------------------------------

Рівняння~\eqref{eq:klein_gordon} описує фонове поле $\bar{\phi}(t)$.
Лінеаризуємо його для малого збурення $\delta\phi = \phi - \bar{\phi}$,
нехтуючи членом $V''(\bar{\phi})\,\delta\phi$ у субхаббловому режимі
($k^2 \gg \mathcal{H}^2$, де масою поля можна знехтувати), і переходимо
до конформного часу $\eta$ та Фур'є-простору. Використовуючи
$\dot{f} = f'/a$ і $\ddot{f} = (f'' - \mathcal{H}f')/a^2$\footnote{%
З означення $d\eta = dt/a$ маємо $\dot{f} = f'/a$.
Диференціюємо ще раз:
$\ddot{f} = \dfrac{d}{dt}\!\left(\dfrac{f'}{a}\right)
= \dfrac{f''}{a^2} - \dfrac{a'}{a^2}\dfrac{f'}{a}
= \dfrac{1}{a^2}(f'' - \mathcal{H}f')$.},
отримуємо для $k$-ї моди \cite{mukhanov_2005}:
\begin{equation}\label{eq:fourier_field}
  \delta\phi_k'' + 2\mathcal{H}\,\delta\phi_k'
  + k^2\,\delta\phi_k = 0,
\end{equation}
де штрих~--- похідна за $\eta$, $k \equiv |\mathbf{k}|$~---
хвильове число у координатах що рухаються з розширенням
(фізична довжина хвилі $\lambda_\mathrm{phys} = a/k$ росте з часом,
тоді як $k$ залишається сталим). Поведінка кожної моди визначається
співвідношенням між $\lambda_\mathrm{phys}$ та горизонтом Хаббла $H^{-1}$:

\begin{itemize}
  \item \textbf{Субхаббловий режим ($k \gg \mathcal{H}$):}
        Для осцилюючої моди $\delta\phi_k' \sim k\,\delta\phi_k$,
        тому хабблівський член за порядком величини:
        $2\mathcal{H}\delta\phi_k' \sim \mathcal{H}k\,\delta\phi_k
        \ll k^2\delta\phi_k$,
        оскільки $\mathcal{H} \ll k$.
        Рівняння~\eqref{eq:fourier_field} зводиться до звичайного
        гармонічного осцилятора:
        \begin{equation*}
          \delta\phi_k'' + k^2\,\delta\phi_k \approx 0.
        \end{equation*}
        Це рівняння вільного поля у пласкому просторі Мінковського~---
        таке саме, як для фотона або будь-якого іншого квантового поля.
        Його квантовий вакуумний стан добре відомий: це стан з мінімальною
        енергією, де кожна мода $k$ знаходиться у нульових коливаннях з
        амплітудою $\sim 1/\sqrt{2k}$ (у одиницях $\hbar = c = 1$).

        У розширюваному просторі де-Сіттера цей стан називається
        \emph{вакуумом Банча--Девіса} \cite{bunch_davies_1978}~---
        це природний вибір початкового стану: на достатньо малих
        масштабах ($k \gg \mathcal{H}$) кривина простору непомітна,
        і поле «не знає», що воно у де-Сіттері. А отже його вакуум
        збігається зі звичайним вакуумом Мінковського~--- тим самим,
        який ми знаємо з КТП у пласкому просторі. Саме цей стан і
        є природною початковою умовою: ми не вигадуємо нічого нового,
        а просто беремо звичайний квантовий вакуум там, де фізика
        ще «не відчула» розширення. З урахуванням
        масштабного фактора $a$ нормування амплітуди:
        \begin{equation*}
          |\delta\phi_k| \sim \frac{1}{a\sqrt{2k}}.
        \end{equation*}
        Множник $1/a$ має простий фізичний зміст: хвильове число
        $k$ визначене у координатах що рухаються разом з розширенням
        (у тих самих $x^i$, де записана метрика FLRW~--- у них
        галактики стоять на місці, а простір між ними розтягується).
        Тому $k$ не змінюється з часом, але фізична довжина хвилі
        $\lambda_\mathrm{phys} = a/k$ розтягується разом з простором.
        Енергія моди «розмазується» по більшому об'єму, і амплітуда
        спадає як $1/a$~--- так само як температура реліктового
        випромінювання спадає як $1/a$ через космологічне
        розтягування фотонів.
  \item \textbf{Суперхаббловий режим ($k \ll \mathcal{H}$):}
        Градієнтний член $k^2\delta\phi_k$ стає нехтовним
        порівняно з $\mathcal{H}^2\delta\phi_k$, і рівняння
        \eqref{eq:fourier_field} перетворюється з осцилятора
        на рівняння загасання:
        \begin{equation*}
          \delta\phi_k'' + 2\mathcal{H}\,\delta\phi_k' \approx 0.
        \end{equation*}
        Член $k^2$ — це «пружина» осцилятора: саме він змушує
        моду коливатися. Коли мода виходить за горизонт,
        $k/\mathcal{H} \to 0$ і пружина «зникає». Фізично:
        два кінці хвилі розлетілися швидше за швидкість світла
        і більше не перебувають у причинному контакті — хвиля
        не може «підтримувати» себе і замерзає.
        Розв'язок: $\delta\phi_k = \mathrm{const}$ (домінуюча мода).
        Мода «виходить за горизонт» і заморожується
        \cite{hawking_1982}.
\end{itemize}

У суперхаббловому ліміті амплітуда виходить на стаціонарне
значення \cite{guth_pi_1982}:
\begin{equation*}
  |\delta\phi_k| \;\to\; \frac{H}{\sqrt{2k^3}}.
\end{equation*}

Звідси первинний безрозмірний спектр потужності флуктуацій поля
\cite{mukhanov_chibisov_1981}:
\begin{equation}\label{eq:field_power_spec}
  \mathcal{P}_{\delta\phi}(k)
    \equiv \frac{k^3}{2\pi^2}|\delta\phi_k|^2
    = \left(\frac{H}{2\pi}\right)^{\!2}.
\end{equation}

\textit{Зауваження: де-Сіттерівська температура і плоский спектр.}
Результат~\eqref{eq:field_power_spec} має глибший зміст, ніж
просто «амплітуда замороженої моди». Простір де-Сіттера має
горизонт подій~--- так само, як чорна діра. Спостерігач у такому
просторі не має доступу до інформації поза горизонтом, і це
обмеження, за Гіббонсом і Гокінгом \cite{gibbons_hawking_1977},
призводить до теплового випромінювання вакуумних флуктуацій з
характерною температурою
\begin{equation*}
  T_\mathrm{dS} = \frac{H}{2\pi}.
\end{equation*}
Саме ця температура визначає амплітуду заморожених флуктуацій~---
вираз $H/(2\pi)$ у формулі~\eqref{eq:field_power_spec} є не
збігом, а прямим наслідком квантової термодинаміки горизонту.

З цього випливає ключова властивість спектра: оскільки $H \approx
\mathrm{const}$ під час інфляції, амплітуда $\mathcal{P}_{\delta\phi}
\propto H^2$ практично не залежить від $k$~--- тобто флуктуації
всіх масштабів народжуються з однаковою «температурою».
Це і є спектр Харісона--Зельдовича, майже інваріантний за
масштабом. Невелике відхилення від точної сталості~---
спектральний індекс $n_s \neq 1$~--- виникає тому, що $H$
повільно спадає під час slow-roll, і детально розглядається
в наступному підрозділі.

%% --------------------------------------------------------
\subsection{Збурення кривини і первинний спектр}
%% --------------------------------------------------------

Збурення $\delta\phi = \phi - \bar\phi$ залежить від того, як проведена
координатна сітка: різні системи координат дають різне значення $\delta\phi$
в одній і тій самій фізичній події.
Тому флуктуації поля $\delta\phi$ самі по собі не мають
калібрувально-інваріантного змісту і не зручні як початкова умова.
Щоб отримати величину, яка однакова в будь-якій калібровці і може
однозначно передаватися від інфляції до подальшої еволюції, вводять
\emph{збурення просторової кривини} $\mathcal{R}_k$, означене
через потенціали $\Phi$ і $\Psi$ у попередньому
розділі~\ref{sec:metric}.

\textit{Чому $\mathcal{R}_k$ є сталим за горизонтом?}
Рівняння руху для $\mathcal{R}_k$ у Фур'є-просторі має вигляд
\cite{bardeen_1980}:
\begin{equation}\label{eq:R_eom}
  \mathcal{R}_k'' + \frac{2z'}{z}\,\mathcal{R}_k' + k^2\mathcal{R}_k = 0,
  \qquad z \equiv \frac{a\dot{\bar\phi}}{H}.
\end{equation}
У суперхаббловому режимі ($k \ll \mathcal{H}$) член $k^2\mathcal{R}_k$
стає нехтовним. Рівняння зводиться до
$\mathcal{R}_k'' + (2z'/z)\,\mathcal{R}_k' = 0$,
розв'язок якого~--- $\mathcal{R}_k = \mathrm{const}$ (зростаюча мода)
і $\mathcal{R}_k \propto \int d\eta/z^2$ (загасаюча мода).
Загасаюча мода спадає і нею нехтують; зростаюча мода заморожується
на сталому значенні. Ця сталість не залежить від деталей реогріву і
постінфляційної еволюції~--- тому $\mathcal{R}_k$ є надійною
«початковою умовою», що передається від інфляції до гарячого
Всесвіту незмінною \cite{bardeen_1980}.

Тут нам достатньо знати, що в конформно-ньютонівській калібровці
$\mathcal{R}_k$ лінійно пов'язане з $\delta\phi_k$ через швидкість
кочення класичного фону~--- це безпосередньо випливає з означення
$\mathcal{R}$ через $\Phi$ і $\Psi$ при підстановці рівнянь руху:
\begin{equation}\label{eq:curvature_perturbation}
  \mathcal{R}_k = -\frac{H}{\dot{\bar{\phi}}}\,\delta\phi_k.
\end{equation}
Підставляючи~\eqref{eq:field_power_spec}, отримуємо первинний
спектр кривини:
\begin{equation}\label{eq:curvature_power_spec}
  \mathcal{P}_{\mathcal{R}}(k)
    = \left(\frac{H}{\dot{\bar{\phi}}}\right)^{\!2}
      \mathcal{P}_{\delta\phi}
    = \frac{H^4}{4\pi^2\dot{\bar{\phi}}^2}.
\end{equation}

Ця формула є центральною для всього подальшого аналізу. У
стандартному slow-roll режимі $\dot{\bar\phi}$ змінюється
повільно, тому $\mathcal{P}_{\mathcal{R}}(k)$ є майже не
залежним від $k$~--- так званий спектр Харісона--Зельдовича.
Відхилення від плоского спектру параметризується спектральним
індексом:
\begin{equation}\label{eq:ns}
  n_s - 1 \equiv \frac{d\ln\mathcal{P}_{\mathcal{R}}}{d\ln k}
            = -2\varepsilon - \eta_V,
\end{equation}
де $\varepsilon$ та $\eta_V$~--- параметри slow-roll з
формули~\eqref{eq:slow_roll_pot}, обчислені в момент виходу
моди $k$ за горизонт.

\textit{Виведення~\eqref{eq:ns}.}
Кожна мода $k$ виходить за горизонт у момент $k = \mathcal{H}
= aH$, тобто число $e$-folds до кінця інфляції $N_k$ залежить
від $k$: $dN = -H\,dt$, звідки $d\ln k = dN$ вздовж траєкторії
slow-roll. Тому:
\begin{equation*}
  n_s - 1
    = \frac{d\ln\mathcal{P}_{\mathcal{R}}}{d\ln k}
    = \frac{d}{dN}\ln\!\left(\frac{H^4}{\dot{\bar\phi}^2}\right)
    = 4\frac{\dot H}{H^2} - 2\frac{\ddot{\bar\phi}}{H\dot{\bar\phi}}
    = -4\varepsilon + 2\eta_\mathrm{H},
\end{equation*}
де використано $\dot H = -\varepsilon H^2$ (з означення~\eqref{eq:slow_roll_def})
та $\eta_\mathrm{H} \equiv -\ddot{\bar\phi}/(H\dot{\bar\phi})$
(знак мінус у означенні).
У потенціальному наближенні slow-roll зв'язок між $\eta_\mathrm{H}$
і $\eta_V$ отримується диференціюванням attractor-рівняння
$3H\dot{\bar\phi} \approx -V'(\bar\phi)$ за часом $t$:
\begin{equation*}
  3\dot H\dot{\bar\phi} + 3H\ddot{\bar\phi} \approx -V''(\bar\phi)\dot{\bar\phi}.
\end{equation*}
Ділимо на $3H^2\dot{\bar\phi}$:
\begin{equation*}
  \frac{\dot H}{H^2} + \frac{\ddot{\bar\phi}}{H\dot{\bar\phi}}
  \approx -\frac{V''}{3H^2}
  \approx -\frac{M_\mathrm{Pl}^2 V''}{V} \cdot \frac{V}{3H^2 M_\mathrm{Pl}^2}
  \approx -\eta_V \cdot 1,
\end{equation*}
де в останньому кроці $V/(3H^2 M_\mathrm{Pl}^2) \approx 1$ у slow-roll.
Отже, з $\dot H/H^2 = -\varepsilon$ і $\eta_\mathrm{H} = -\ddot{\bar\phi}/(H\dot{\bar\phi})$:
\begin{equation*}
  -\varepsilon - \eta_\mathrm{H} \approx -\eta_V
  \quad\Rightarrow\quad
  \eta_\mathrm{H} \approx \varepsilon - \eta_V.
\end{equation*}
Точніший рахунок з урахуванням поправок порядку $\varepsilon^2$
дає \cite{mukhanov_2005}:
\[
  \eta_\mathrm{H} \approx \varepsilon - \tfrac{1}{2}\eta_V,
\]
звідки:
\begin{equation*}
  n_s - 1 = -4\varepsilon + 2\bigl(\varepsilon -
  \tfrac{1}{2}\eta_V\bigr) = -2\varepsilon - \eta_V. \qquad\square
\end{equation*}
Planck~2018 вимірює $n_s = 0.9649 \pm 0.0042$~--- ледь нахилений,
але не плоский спектр, що є прямим підтвердженням slow-roll.

Проте якщо на шляху інфлатона є локальна аномалія потенціалу
(сходинка, ділянка USR), швидкість $\dot{\bar\phi}$ різко падає.
Оскільки $\mathcal{P}_{\mathcal{R}} \propto \dot{\bar\phi}^{-2}$,
флуктуації, що виходять за горизонт у цей момент, зазнають
локального підсилення~--- формується \textbf{зламаний спектр}.
Механізм і наслідки детально розглядаються у
§\,\ref{sec:broken}.

%% --------------------------------------------------------
\subsection{Заморожування та розморожування мод}
%% --------------------------------------------------------

Еволюція первинних збурень підпорядковується чіткій часовій
схемі:
\begin{enumerate}
  \item \textbf{Заморожування (інфляція):} Простір розширюється
        прискорено ($a \propto e^{Ht}$), горизонт Хаббла
        $H^{-1}$ залишається майже сталим. Масштаби збурень
        послідовно покидають горизонт ($\lambda_\mathrm{phys} >
        H^{-1}$), і збурення $\mathcal{R}_k$ заморожуються,
        зберігаючи відбиток локальної мікрофізики $V(\phi)$
        \cite{bardeen_1980}.
  \item \textbf{Розморожування (постінфляційні епохи):}
        Розширення сповільнюється, горизонт Хаббла росте
        швидше за $a$. Збурення послідовно повертаються в
        горизонт~--- спочатку малі масштаби (великі $k$),
        пізніше великі (малі $k$) \cite{mukhanov_2005}.
\end{enumerate}

Повернувшись в горизонт, кожна мода $\mathcal{R}_k$ стає
початковою умовою для класичної гідродинаміки фотон-баріонної
плазми. Але перш ніж плазма виникне~--- потрібен реогрів:
перетворення енергії інфлатона на гарячі частинки Стандартної
моделі. Саме з цього починається наступний розділ.


%% --------------------------------------------------------
\section{Реогрів і народження матерії}
\label{sec:reheating}
%% --------------------------------------------------------

Коли інфлатон досягає мінімуму потенціалу $V(\phi)$, умова
інфляції порушується ($\varepsilon \to 1$) і прискорене
розширення припиняється. Але Всесвіт у цей момент
перебуває у парадоксальному стані: він розширився в
$e^{50}$--$e^{60}$ разів, охолов майже до абсолютного нуля
і є практично порожнім~--- вся енергія зосереджена у
вигляді однорідних класичних коливань поля $\phi$ навколо
мінімуму. Жодної плазми, жодних часток Стандартної моделі.

Перехід від цього «холодного» стану до гарячого
ультрарелятивістського Всесвіту, з якого починається
стандартна космологія Великого вибуху, описується теорією
\emph{реогріву} \cite{linde_1982, albrecht_1982}. Динаміку
реогріву поділяють на дві послідовні стадії: спочатку
нелінійний вибуховий \textbf{пререогрів}, де параметричний
резонанс за лічені осциляції перекачує енергію з поля у
частинки, а потім повільний \textbf{пертурбативний розпад},
де залишки інфлатона розпадаються і система термалізується.

%% --------------------------------------------------------
\subsection{Пререогрів та параметричний резонанс}
%% --------------------------------------------------------

Поблизу мінімуму параболічного потенціалу
$V(\phi) = \frac{1}{2}m^2\phi^2$ інфлатон осцилює як
$\phi(t) \approx \Phi_0(t)\sin(mt)$, де амплітуда
$\Phi_0(t)$ повільно затухає через хабблівське тертя.
Якщо інфлатон пов'язаний з іншим скалярним полем $\chi$
через член взаємодії $g^2\phi^2\chi^2$, то ефективна маса
поля $\chi$ осцилює в часі: $m_\chi^2(t) \supset
g^2\Phi_0^2\sin^2(mt)$. Рівняння руху для $k$-ї моди
$\chi_k$ набуває форми рівняння Матьє
\cite{kofman_1994, kofman_1997}:
\begin{equation}\label{eq:mathieu}
  \ddot{\chi}_k + 3H\dot{\chi}_k
  + \bigl[k_\mathrm{phys}^2 + m_\chi^2
         + g^2\Phi_0^2\sin^2(mt)\bigr]\chi_k = 0.
\end{equation}

Рівняння з періодично змінними коефіцієнтами має зони
нестабільності, де амплітуда розв'язку зростає
експоненціально~--- \emph{параметричний резонанс}.
Народження часток $\chi$ відбувається не поштучно, а
вибухово: бозонні поля мають стимульоване народження у
вже заповнені стани (аналог ефекту Бозе--Айнштайна), тому
заселеність кожної резонансної моди зростає за лічені
осциляції інфлатона. Цей процес~--- \textbf{пререогрів}
(preheating)~--- за короткий час перекачує значну частку
енергії поля у частинки $\chi$.

%% --------------------------------------------------------
\subsection{Пертурбативний розпад та термалізація}
%% --------------------------------------------------------

Коли амплітуда $\Phi_0$ падає нижче порогу резонансу,
параметричний резонанс вимикається. Подальший розпад
залишків інфлатона відбувається повільно і описується
пертурбативною теорією: у рівняння збереження енергії
вводиться феноменологічна ширина розпаду $\Gamma_\phi$:
\begin{align}
  \dot{\rho}_\phi + 3H\rho_\phi
    &= -\Gamma_\phi\,\rho_\phi,
    \label{eq:reheat_rho_phi} \\
  \dot{\rho}_r + 4H\rho_r
    &= +\Gamma_\phi\,\rho_\phi.
    \label{eq:reheat_rho_r}
\end{align}
Перше рівняння~--- розпад інфлатона; друге~--- накопичення
радіації. Коефіцієнти $3H$ і $4H$ відображають розбавлення
густин розширенням: $\rho_\phi \propto a^{-3}$ (матерія),
$\rho_r \propto a^{-4}$ (радіація).

Поки $H \gg \Gamma_\phi$, розширення домінує і розпад
відбувається повільно. Процес завершується при
$H \sim \Gamma_\phi$: залишки інфлатона розпадаються майже
миттєво, а релятивістські продукти термалізуються через
пружне і непружне розсіяння, встановлюючи єдину температуру.

\begin{keybox}
\textbf{Температура реогріву $T_\mathrm{rh}$:}
визначається з умови $H^2 = \frac{8\pi G}{3}\rho_r$ при
$H \approx \Gamma_\phi$ і рівноважного розподілу радіації
$\rho_r = \frac{\pi^2}{30}g_* T^4$:
\begin{equation}\label{eq:t_reheating}
  T_\mathrm{rh}
    = \left(\frac{90}{8\pi^3 g_*}\right)^{\!1/4}
      \sqrt{M_\mathrm{Pl}\,\Gamma_\phi},
\end{equation}
де $g_* \sim 100$~--- ефективна кількість
релятивістських ступенів вільності Стандартної моделі.
\end{keybox}

Температура $T_\mathrm{rh}$ задає початкові умови для
всієї подальшої термічної історії: баріогенезису, фазових
переходів і утворення темної матерії. Для нашої задачі
важливо інше: саме в момент термалізації гаряча
фотон-баріонна плазма «успадковує» просторові
неоднорідності кривини $\mathcal{R}_k$, заморожені під час
інфляції. З цього моменту вони більше не є статичними~---
тиск фотонного газу змушує плазму коливатися, і саме ці
акустичні коливання ми бачимо сьогодні у вигляді
температурних плям на карті CMB (рис.~\ref{fig:planck_map})
і піків на її кутовому спектрі (рис.~\ref{fig:cmb_tt}).

\begin{figure}[htbp]
\centering
\includegraphics[width=\textwidth]{planck_cmb_map.png}
\caption{Карта температурних анізотропій CMB від місії
         Planck~2018 \cite{planck_2018}. Кольори відображають
         відхилення від середнього $T_0 = 2.725\,\mathrm{K}$
         у діапазоні $\pm 300\,\mu\mathrm{K}$~--- спостережуваний
         відбиток первинних флуктуацій інфлатона після
         ${\sim}380\,000$ років еволюції.
         Зображення: ESA and the Planck Collaboration
         \cite{esa_planck_cmb_image}.}
\label{fig:planck_map}
\end{figure}

%% --------------------------------------------------------
\section{Гідродинаміка фотон-баріонної плазми та акустичні осциляції}
\label{sec:bao}
%% --------------------------------------------------------

Після реогріву Всесвіт входить у радіаційно-доміновану епоху.
Температура настільки висока, що баріонна матерія повністю
іонізована: протони та електрони утворюють плазму, в якій
фотони інтенсивно розсіюються на вільних електронах
(томсонівське розсіяння). Середня довжина вільного пробігу
фотона набагато менша за космологічний горизонт, тому
випромінювання і баріони поводяться як єдина
\emph{фотон-баріонна плазма} \cite{mukhanov_2005}.

Первинні неоднорідності кривини $\mathcal{R}_k$, успадковані
від інфляції, стають початковими умовами для цієї плазми.
Подальша динаміка визначається конкуренцією двох сил:
гравітаційне притягання намагається стиснути речовину у
потенціальні ями $\Phi$, а внутрішній тиск фотонного газу
чинить опір. Результат~--- акустичні коливання.

%% --------------------------------------------------------
\subsection{Рівняння акустичного осцилятора}
%% --------------------------------------------------------

Щоб описати ці коливання кількісно, потрібно знати
швидкість звуку в плазмі.
Швидкість звуку визначається як $c_s^2 = \partial p/\partial\rho$.
Оскільки тиск цілком фотонний ($p_\gamma = \bar{\rho}_\gamma/3$,
а баріонний тиск нехтовно малий~--- нерелятивістський протонний
газ при $T \sim 10^4$~К має $c_s^{(b)} \sim \sqrt{T/m_p}
\sim 10^{-5}\,c$), то $c_s = 1/\sqrt{3} \approx 0.577\,c$
у чистій фотонній плазмі.

В релятивістській гідродинаміці маса спокою --- не єдиний
внесок в інерцію: тиск теж чинить опір прискоренню, і
роль \enquote{ефективної маси} відіграє ентальпія $\rho + p$.
Для фотонів $\rho_\gamma + p_\gamma = \tfrac{4}{3}\rho_\gamma$~---
саме це число визначає, наскільки «важкою» є фотонна компонента
як партнер баріонів у спільних коливаннях. Зручно ввести
безрозмірний параметр, що показує відносну вагу баріонної
інерції порівняно з фотонною (baryon loading parameter):
\begin{equation*}
  R(\eta) \equiv
    \frac{\bar{\rho}_b}{\rho_\gamma + p_\gamma}
    = \frac{3\bar{\rho}_b(\eta)}{4\bar{\rho}_\gamma(\eta)}
    \propto a(\eta).
\end{equation*}
$R$ зростає з часом, бо $\bar{\rho}_b \propto a^{-3}$
спадає повільніше за $\bar{\rho}_\gamma \propto a^{-4}$.
Фізична аналогія: додаємо масу до маятника без зміни
пружини~--- частота коливань падає.
Підставляючи $R$ у вираз для швидкості звуку:
\begin{equation}\label{eq:sound_speed}
  c_s(\eta) = \frac{1}{\sqrt{3\bigl(1 + R(\eta)\bigr)}}.
\end{equation}
При $R \to 0$ отримуємо $c_s \to 1/\sqrt{3}$~--- чиста
фотонна плазма; при $R \gg 1$ (пізня епоха, багато баріонів)
$c_s \to 0$.

Акустичні коливання фізично можливі лише на субхабблових
масштабах ($k\eta \gg 1$): тільки тоді різні точки плазми
перебувають у причинному контакті і можуть коливатися
узгоджено.

Виведення рівняння осцилятора починається з двох рівнянь
гідродинаміки у Фур'є-просторі для збурення фотонної густини
$\delta_\gamma \equiv \delta\rho_\gamma/\bar{\rho}_\gamma$
та дивергенції поля швидкостей плазми $\theta \equiv \nabla\cdot\mathbf{v}$
\cite{dodelson_2020}:

\textbf{Рівняння неперервності} (збереження фотонів):
\begin{equation*}
  \delta_\gamma' = -\frac{4}{3}\,\theta + 4\Phi',
\end{equation*}
де $4/3$ виникає з рівняння стану $p_\gamma = \rho_\gamma/3$,
а $\Phi'$ — через зміну об'єму в збуреній метриці.

\textbf{Рівняння Ейлера} (баланс тиску і гравітації) для
об'єднаної фотон-баріонної рідини з ефективною інерцією $(1+R)$:
\begin{equation*}
  \theta' + \mathcal{H}\frac{R}{1+R}\,\theta
  = \frac{k^2}{4(1+R)}\,\delta_\gamma + k^2\Psi.
\end{equation*}
Член $\mathcal{H}R/(1+R)$ — це хабблівське тертя, ефективне
лише для баріонної компоненти ($R$); фотони не відчувають тертя.

Диференціюємо рівняння неперервності за $\eta$ і підставляємо
$\theta'$ з рівняння Ейлера. У конформно-ньютонівському
калібруванні $\Phi = \Psi$ і на субхабблових масштабах
членами з $\mathcal{H}$ нехтуємо. Після скорочень отримуємо
рівняння вимушеного осцилятора \cite{dodelson_2020}:
\begin{equation}\label{eq:acoustic_oscillator}
  \frac{d}{d\eta}\bigl[(1+R)\,\delta_\gamma'\bigr]
  + \frac{k^2}{3}\,\delta_\gamma
  = -\frac{k^2}{3}(1+R)\,\Phi.
\end{equation}
Ліва частина~--- осцилятор зі змінною масою $(1+R)$
і частотою $\omega = kc_s$; коефіцієнт $1/3$ є саме
$c_s^2 = 1/(3(1+R))$ з~\eqref{eq:sound_speed}. Права
частина~--- гравітаційне вимушення потенціалом $\Phi$,
що зміщує положення рівноваги. Кожна мода $k$ коливається
незалежно, а оскільки всі вони стартують з однакової
початкової фази $\mathcal{R}_k$~--- на відміну від
некогерентних джерел типу топологічних дефектів~--- їхні
осциляції є \emph{когерентними}: моди з однаковим $k$
перебувають в одній фазі незалежно від просторового
положення. Саме когерентність породжує чіткі акустичні
піки на спектрі CMB.

%% --------------------------------------------------------
\subsection{Рекомбінація, звуковий горизонт і BAO}
%% --------------------------------------------------------

Акустичні коливання тривають до \emph{рекомбінації}
($z_\mathrm{rec} \approx 1100$, $T \approx 3000$~К):
протони захоплюють електрони, іонізація припиняється,
томсонівське розсіяння вимикається і фотони починають
вільно поширюватися. Плазма «заморожується» в тому стані,
в якому її застала рекомбінація.

Максимальна відстань, яку встигла пройти звукова хвиля
від початку гарячої епохи до рекомбінації,~---
\emph{радіус звукового горизонту}:
\begin{equation}\label{eq:sound_horizon}
  r_s(\eta_\mathrm{rec})
    = \int_0^{\eta_\mathrm{rec}} c_s(\eta)\,d\eta
    \approx 147\,\mathrm{Mpc}.
\end{equation}
Це єдиний лінійний масштаб у задачі. Моди, що в момент
рекомбінації перебували у точці максимального стиснення
або розширення, проявляються у температурному спектрі
CMB як серія \emph{акустичних піків}~--- докладніше
у §\,\ref{sec:sachs}.

Але акустичні хвилі залишають слід не лише у фотонах.
Після рекомбінації фотони пішли геть, а баріонна речовина
«замерзла» у вигляді сферичної оболонки радіусом $r_s$
навколо кожної початкової неоднорідності темної матерії
(рис.~\ref{fig:bao_profile}). Це~--- \textbf{баріонні
акустичні осциляції (BAO)} \cite{eisenstein_2005}.

\begin{figure}[b!]
\centering
\begin{tikzpicture}
  \begin{axis}[
    width=\linewidth,
    height=0.7\linewidth,
    xmin=0, xmax=250,
    ymin=-0.05, ymax=1.1,
    xlabel={Відстань від первинного збурення, $r$ [Мпк]},
    ylabel={Контраст густини баріонів,
            $\Delta\rho_b/\bar{\rho}_b$},
    grid=both,
    grid style={line width=.1pt, color=gray!20},
    major grid style={line width=.2pt, color=gray!40},
    tick label style={font=\footnotesize},
    label style={font=\small},
    legend style={at={(0.95,0.95)}, anchor=north east,
                  font=\footnotesize, cells={anchor=west}},
    axis lines=left,
    clip=false
  ]
    \addplot[
      domain=0:240, samples=200,
      color=accent, line width=1.5pt
    ] {exp(-0.015*x) + 0.35*exp(-((x-147)/12)^2)};
    \addlegendentry{Профіль густини речовини}

    \draw[dashed, color=theoryred, line width=1pt]
      (axis cs:147,-0.05) -- (axis cs:147,0.9);
    \node[anchor=south, color=theoryred, font=\footnotesize]
      at (axis cs:147,0.9) {$r_s \approx 147\text{ Мпк}$};

    \node[anchor=west, color=gray, font=\scriptsize]
      at (axis cs:165,0.2) {Затухання Шовка};
    \draw[->, color=gray, >=stealth]
      (axis cs:162,0.18) -- (axis cs:155,0.1);
  \end{axis}
\end{tikzpicture}
\caption{Схематичний профіль поширення первинного
  адіабатичного збурення у баріонній компоненті плазми
  (модельна крива). Надлишок речовини на радіусі
  звукового горизонту $r_s$ відповідає сучасному
  масштабу BAO.}
\label{fig:bao_profile}
\end{figure}

У сучасну епоху масштаб BAO проявляється як фіксований
надлишок у просторовій кореляційній функції розподілу
галактик на відстані $\sim 147$~Мпк. Оскільки $r_s$
жорстко визначений лінійною фізикою раннього
Всесвіту~\eqref{eq:sound_horizon}, BAO~--- ідеальна
«стандартна лінійка»: вимірюючи цей масштаб на різних
червоних зміщеннях, можна реконструювати $H(z)$ і
параметри темної енергії незалежно від CMB.

Той самий звуковий горизонт $r_s$, що визначає масштаб
BAO у розподілі галактик, відповідає і за положення
акустичних піків у температурному спектрі CMB. Як саме
ці піки формуються~--- і які інші ефекти на них
накладаються~--- розглядається у наступному розділі.


%% --------------------------------------------------------
\section{Ефект Закса--Вольфа та формування спектра CMB}
\label{sec:sachs}
%% --------------------------------------------------------

Кутова анізотропія температури реліктового випромінювання (CMB), що спостерігається нами сьогодні, є безпосереднім відбитком («знімком») стану акустичних осциляцій фотон-баріонної плазми в епоху рекомбінації Всесвіту ($z_{\mathrm{rec}} \approx 1100$). Флуктуації температури у заданому напрямку на небі $\hat{\mathbf{n}}$ розкладаються за сферичними гармоніками, а їхні статистичні властивості повністю описуються кутовим спектром потужності $\mathcal{D}_\ell^{TT} \equiv \frac{\ell(\ell+1)}{2\pi}C_\ell^{TT}$, де мультиполь $\ell$ обернено пропорційний кутовому масштабу $\theta \sim 180^\circ / \ell$ \cite{dodelson_2020}.

Спостережувані варіації температури $\delta T / T$ формуються внаслідок дії трьох основних фізичних механізмів на поверхні останнього розсіяння (LSS) \cite{hu_sugiyama_1995}:
\begin{equation}\label{eq:cmb_sources}
  \frac{\delta T}{\text{T}}(\hat{\mathbf{n}}) = \left( \frac{1}{4}\delta_\gamma + \Phi \right)_{\mathrm{LSS}} + \left( \hat{\mathbf{n}} \cdot \mathbf{v}_b \right)_{\mathrm{LSS}} + \int_{\eta_{\mathrm{LSS}}}^{\eta_0} \left( \Phi' + \Psi' \right) d\eta.
\end{equation}

Кожен із доданків у виразі (\ref{eq:cmb_sources}) відповідає за свій масштаб на спектрі:
\begin{enumerate}
  \item \textbf{Звичайний ефект Закса--Вольфа ($\frac{1}{4}\delta_\gamma + \Phi$):}
        Домінує на великих кутових масштабах (малі мультиполі $\ell \lesssim 100$),
        які на момент рекомбінації ще перебували поза горизонтом Хаббла.
        Він поєднує в собі власну термодинамічну анізотропію фотонів та
        гравітаційне червоне зміщення, якого фотони зазнають при виході
        з потенціальних ям $\Phi$.

        \textit{Адіабатичні збурення і умова $\frac{1}{4}\delta_\gamma = -\frac{2}{3}\Phi$.}
        Збурення називаються \emph{адіабатичними}, якщо всі компоненти
        (фотони, баріони, темна матерія) збурені узгоджено: локальний
        стан плазми відрізняється від середнього лише температурою, а
        не складом. Формально це означає, що відношення числових густин
        зберігається: $\delta(n_b/n_\gamma) = 0$. Оскільки
        $n_\gamma \propto T^3$ і $n_b \propto a^{-3}$, умова адіабатичності
        зводиться до $\delta_\gamma = \frac{4}{3}\delta_b$.

        На суперхабблових масштабах у радіаційно-домінованій епосі
        рівняння збереження для фотонів і метричний потенціал $\Phi$
        пов'язані через рівняння Пуассона:
        $k^2\Phi \approx -4\pi G a^2 \bar\rho\,\delta$.
        У граничному випадку $k \to 0$ (масштаби значно більші за горизонт)
        розв'язок для адіабатичних початкових умов дає
        \cite{hu_sugiyama_1995}:
        \begin{equation*}
          \frac{1}{4}\delta_\gamma = -\frac{2}{3}\Phi,
        \end{equation*}
        де мінус означає, що надлишок густини фотонів ($\delta_\gamma > 0$)
        відповідає потенціальній ямі ($\Phi < 0$): плазма стиснута саме
        там, де є гравітаційний мінімум. Підставляючи у вираз для ефекту
        Закса--Вольфа:
        \begin{equation*}
          \frac{1}{4}\delta_\gamma + \Phi
          = -\frac{2}{3}\Phi + \Phi = \frac{1}{3}\Phi.
        \end{equation*}
        Знак «$+$»: фотони виходять з ями ($\Phi < 0$) і зазнають
        гравітаційного червоного зміщення, тобто виглядають
        \emph{холоднішими}. Ефект $\frac{1}{3}\Phi < 0$ означає
        дефіцит температури у напрямку потенціальної ями~---
        саме це формує плато Закса--Вольфа на великих масштабах.
  \item \textbf{Доплерівський ефект ($\hat{\mathbf{n}} \cdot \mathbf{v}_b$):} Зумовлений кінематичним рухом (швидкістю $\mathbf{v}_b$) фотон-баріонної плазми у полі збурень. Цей внесок зміщений за фазою відносно ефекту Закса--Вольфа і частково заповнює мінімуми між акустичними піками.
  \item \textbf{Інтегральний ефект Закса--Вольфа (ISW, інтегральний член):} Виникає, коли фотони CMB по дорозі до спостерігача пролітають крізь гравітаційні потенціали, що змінюються у часі ($\Phi' \neq 0$). Ранній ISW проявляється відразу після рекомбінації через залишковий вплив радіації, а пізній ISW~--- на надвеликих масштабах ($\ell < 10$) в сучасну епоху через прискорене розширення Всесвіту під дією темної енергії.
\end{enumerate}

%% --------------------------------------------------------
\subsection{Акустичні піки та затухання Шовка}
%% --------------------------------------------------------

На субхабблових масштабах ($\ell \gtrsim 100$) когерентні звукові коливання формують чітку серію \emph{акустичних піків} (рис.~\ref{fig:cmb_tt}). Перший і найвищий пік при $\ell \approx 220$ ($\theta \sim 1^\circ$) відповідає моді, яка встигла зазнати рівно одного максимального стиснення у потенціальній ямі до моменту рекомбінації. Його положення є чутливим індикатором просторової геометрії Всесвіту \cite{planck_2018}. На мультиполях $\ell \gtrsim 1000$ спектр експоненціально загасає через ефект \emph{дифузійного затухання} (затухання Шовка)~--- фотони мають скінченну довжину вільного пробігу і в процесі дифузії замивають дрібномасштабні флуктуації температури. Характерний масштаб затухання $\ell_D \sim 1000$--$1500$ визначається довжиною дифузії фотонів до рекомбінації \cite{silk_1968}.

\begin{figure}[b!]
\centering
\begin{tikzpicture}
\begin{semilogxaxis}[
    width  = 14.5cm,
    height = 8.5cm,
    xlabel = {Мультипольний момент $\ell$},
    ylabel = {$\mathcal{D}_\ell^{TT}\ [\mu\mathrm{K}^2]$},
    xmin   = 2,   xmax = 2600,
    ymin   = 0,   ymax = 6800,
    xtick  = {2, 10, 30, 100, 300, 1000, 2500},
    xticklabels = {2, 10, 30, 100, 300, 1000, 2500},
    minor xtick = {5,20,50,200,500,2000},
    ytick  = {0, 1000, 2000, 3000, 4000, 5000, 6000},
    grid   = both,
    grid style = {line width=0.3pt, draw=gray!25},
    major grid style = {line width=0.5pt, draw=gray!40},
    tick label style = {font=\small},
    label style      = {font=\small},
    legend style = {
        at={(0.6,0.97)}, anchor=north east,
        font=\small, fill=white, fill opacity=0.9,
        draw=gray!60, inner sep=4pt, row sep=1pt
    },
    clip = true,
]

%--- Theory curve (LCDM, CAMB) ---
\addplot[
    color     = theoryred,
    line width = 1.6pt,
    smooth,
] table[x=ell, y=Dl] {\theorydata};
\addlegendentry{$\Lambda$CDM (CAMB, Planck\,2018 best-fit)}

%--- Planck 2018 data with error bars ---
\addplot[
    only marks,
    mark       = *,
    mark size  = 1.4pt,
    color      = planckblue,
    error bars/.cd,
        y dir  = both,
        y explicit,
] table[
    x    = ell,
    y    = Dl,
    y error minus = errm,
    y error plus  = errp,
] {\planckdata};
\addlegendentry{Planck\,2018 (збіновані дані TT)}

%--- Annotate acoustic peaks ---
\node[font=\tiny, text=gray!70, above] at (axis cs:220,5732) {1-й пік};
\node[font=\tiny, text=gray!70, above] at (axis cs:540,2592) {2-й пік};
\node[font=\tiny, text=gray!70, above] at (axis cs:810,2540) {3-й пік};

%--- Sachs-Wolfe plateau annotation ---
\draw[->, gray!60, thin] (axis cs:6,1300) -- (axis cs:30,1070)
    node[midway, left, font=\tiny, text=gray!60] {плато Закса--Вольфа};

\end{semilogxaxis}
\end{tikzpicture}
\caption{Кутовий TT-спектр потужності CMB. Точки~--- дані Planck~2018
         \cite{planck_2018}.
         Крива~--- базова модель $\Lambda$CDM, розрахована CAMB
         з параметрами місії Planck~2018 \cite{planck_2018}
         ($H_0=67.36$, $\Omega_b h^2=0.02237$, $\Omega_c h^2=0.1200$,
         $\tau=0.0544$, $A_s=2.1\times10^{-9}$, $n_s=0.9649$).
         Видно три акустичні піки та плато Закса--Вольфа на великих
         кутових масштабах ($\ell \lesssim 30$).}
\label{fig:cmb_tt}
\end{figure}


%% --------------------------------------------------------
\section{Поляризація CMB}
\label{sec:pol}
%% --------------------------------------------------------

Температурний спектр $C_\ell^{TT}$, розглянутий у попередньому
розділі, несе інформацію переважно про скалярні збурення.
Але CMB містить і другий канал інформації~--- \emph{поляризацію}.
Вона виникає при томсонівському розсіянні фотонів і дозволяє
розрізнити два типи первинних збурень: скалярні (флуктуації
густини) і тензорні (гравітаційні хвилі). Детекція тензорного
сигналу у поляризації CMB стала б прямим підтвердженням
квантової природи гравітації~--- саме тому пошук B-мод є
однією з центральних задач сучасної космології.

%% --------------------------------------------------------
\subsection*{Механізм поляризації: розсіяння Томсона}
%% --------------------------------------------------------

Фотони поляризуються при останньому розсіянні на електронах.
Томсонівське розсіяння поляризує фотон \emph{лише} якщо навколо електрона
існує \textbf{квадрупольна} температурна анізотропія: якщо температура
однакова з усіх боків~--- поляризації немає.

Безпосередньо до рекомбінації фотонів і електронів ще достатньо для
ізотропізації, тому квадруполь малий і поляризація слабка.
Але під час самої рекомбінації ($\Delta z \sim 80$) оптична глибина
падає швидко~--- квадруполь встигає вирости і поляризація «відбивається»
в поверхні останнього розсіяння.

Напрямок поляризації залежить від \textbf{поперечного} перепаду температури
навколо точки розсіяння.

%% --------------------------------------------------------
\subsection*{E- і B-моди}
%% --------------------------------------------------------

Поляризаційне поле на сфері, як і будь-яке тензорне поле,
однозначно розкладається на дві компоненти~--- аналогічно
до розкладу вектора на потенціальну і вихрову частини.
Ці дві компоненти мають різне фізичне джерело і різний
статус спостережень:

\begin{center}
\begin{tblr}{XXX}
\toprule
Тип & Джерело & Статус \\
\midrule
E-моди (градієнтний, симетричний візерунок) & Скалярні флуктуації & Виміряні \\
B-моди (ротаційний, закручений візерунок)  & Тензорні флуктуації (GW) & Поки не знайдені \\
\bottomrule
\end{tblr}
\end{center}

Чому два типи збурень дають різні моди поляризації~---
питання симетрії, а не лише статистики:

\begin{keybox}
\textbf{Чому скалярні збурення не дають B-мод?}
Скалярні збурення мають потенціальний (паритетно-парний) характер:
$\Phi(\mathbf{x})$ визначає поле швидкостей $\mathbf{v} = \nabla\Phi$.
E-моди мають ту ж симетрію паритету~--- вони парні.
B-моди псевдоскалярні (непарні), і лінійне скалярне збурення їх не генерує.
Тензорне збурення $h_{ij}$ має обидва стани, в тому числі непарний~---
звідси B-моди від гравітаційних хвиль.
\textit{Практичний забруднювач:} гравітаційне лінзування на великих масштабах перетворює
E-моди в B-моди~--- це треба враховувати при пошуку первинних B-мод.
\end{keybox}

Знайти B-моди~--- означає отримати прямий доказ первинних гравітаційних хвиль
і квантової природи гравітації в режимі малих збурень.
Поточна верхня межа: $r < 0.036$ (95\% CL, BICEP/Keck~2021 \cite{bicep_keck_2021}).

Таким чином, CMB є повноцінним «двоканальним» джерелом інформації:
температурний спектр обмежує скалярні збурення і параметри $\Lambda$CDM,
а поляризаційні B-моди~--- єдиний прямий зонд тензорних збурень і
енергетичного масштабу інфляції. Обидва канали разом утворюють
найпотужніший інструмент для перевірки інфляційних моделей~---
включно з тією, що розглядається у наступних розділах.

%% --------------------------------------------------------
\section{Хаотична інфляція Лінде}
\label{sec:chaotic}
%% --------------------------------------------------------

Лінде~\cite{linde_1983} запропонував радикально простий погляд
на початкові умови: значення $\phi$ у різних точках раннього
Всесвіту є хаотичним~--- немає потреби в особливому «стартовому»
стані. Найпростіший потенціал:
\begin{equation*}
  V(\phi) = \tfrac{1}{2}m^2\phi^2.
\end{equation*}

Де $\phi$ велике і $\varepsilon \ll 1$~--- там інфляція іде довго
і виникає величезна причинно зв'язана «бульбашка». Де $\phi$ мале~---
інфляція швидко завершується. Наш спостережуваний Всесвіт~---
одна з таких бульбашок.

\textbf{Структура всередині однієї бульбашки.} Коли інфляція
в нашій бульбашці завершилась, простір всередині неї вже був
практично однорідним~--- розширення загладило всі початкові
неоднорідності $\phi$. Але «практично» не означає «абсолютно»:
залишились крихітні квантові флуктуації, заморожені під час
інфляції. Саме вони стали насінням галактик і є тим, що ми
бачимо на карті CMB.

Всередині однієї бульбашки існує єдиний причинно зв'язаний
простір-час~--- один Всесвіт. Але його розмір після $N \sim 60$
$e$-folds становить $\sim e^{60}$ разів більше за наш
спостережуваний горизонт ($\sim 46$~Гпк). Тому навіть у
межах «нашої бульбашки» існують області, з якими ми ніколи
не матимемо причинного контакту~--- вони назавжди за
нашим горизонтом подій.

\textbf{Між бульбашками.} Простір між різними бульбашками
продовжує розширюватись інфляційно~--- швидше за світло.
Дві бульбашки, що утворились незалежно, розділені областю
вічної інфляції і принципово не можуть обмінятись сигналом.
Це не просто «дуже далеко»~--- між ними немає спільного
майбутнього світлового конуса. У строгому сенсі загальної теорії відносності вони
належать різним причинно відокремленим регіонам
одного глобального простору-часу, а не різним «всесвітам»~---
хоча в популярній літературі їх так і називають.

\textbf{Вічна інфляція.} Квантові флуктуації інфлатона на
суперхабблових масштабах можуть локально \emph{збільшити} $\phi$,
повернувши його назад на схил, якщо флуктуація достатньо велика.
Класичний зсув поля за один $e$-fold: $|\delta\phi_{\rm class}|
\sim |\dot\phi|/H \approx \sqrt{2\varepsilon}\,M_{\rm Pl}$.
Квантова флуктуація за той самий час: $\delta\phi_{\rm quant}
\sim H/(2\pi)$. Умова «самовідтворення»~---
квантова флуктуація перевищує класичне зміщення:
\begin{equation*}
  \frac{H}{2\pi} \gtrsim \sqrt{2\varepsilon}\,M_{\rm Pl}
  \quad\Leftrightarrow\quad
  H^2 \gtrsim M_{\rm Pl}^2
  \quad\Leftrightarrow\quad
  V(\phi) \gtrsim M_{\rm Pl}^4,
\end{equation*}
де в останньому переході використано $H^2 \approx V/(3M_{\rm Pl}^2)$
і $\varepsilon \lesssim \mathcal{O}(1)$.
У таких областях інфляція ніколи не зупиняється глобально:
поки одна бульбашка завершує інфляцію і термалізується,
сусідні регіони продовжують розширюватись, породжуючи нові
бульбашки. Це \textbf{вічна інфляція}~--- мультивсесвіт
як математичний наслідок, а не постулат.

\textbf{Статус.} Потенціал $\tfrac{1}{2}m^2\phi^2$ передбачає
$r \approx 8/N \approx 0.13$~--- вище поточної межі BICEP/Keck
(табл.~\ref{tab:inflation_models}), тому ця конкретна модель
виключена. Але ідея хаотичних початкових умов і вічної інфляції
залишається живою у більш пологих потенціалах (Старобінський,
хіггсова інфляція).

Разом з тим жодна з «успішних» моделей~--- тих, що проходять
обмеження Planck на $n_s$ і $r$~--- не передбачає нічого
особливого на малих масштабах $k \gtrsim 1\,h/$Мпк. Саме там
первинний спектр залишається практично необмеженим
спостереженнями. І саме там у 2022--2023 роках телескоп JWST
зафіксував аномалії, які стандартний спектр пояснити не може.


%% --------------------------------------------------------
\section{Наша задача: зламаний спектр і JWST}
\label{sec:broken}
%% --------------------------------------------------------

%% --------------------------------------------------------
\subsection{Мотивація}
%% --------------------------------------------------------

JWST знайшов масивні галактики при $z \sim 10$--$12$, яких у стандартній
моделі бути не повинно~--- не вистачало часу утворитись. Конкретно:

\begin{itemize}
  \item Labb\'e et al.\ \cite{labbe_2023} виявили шість галактик-кандидатів
        при $z \sim 7$--$10$ із зоряними масами $M_\star \sim 10^{10}$--$10^{11}
        M_\odot$. Сумарна щільність зоряної маси при $z > 7.4$
        перевищує передбачення $\Lambda$CDM у $\sim\!20$ разів і
        виходить за межу барионної маси доступних гало у стандартній моделі.
  \item Finkelstein et al.\ \cite{finkelstein_2023} знайшли яскраву галактику
        Maisie's Galaxy при $z \approx 11.4$~--- за ${\sim}390$ млн років
        після Великого Вибуху~--- з темпом зореутворення на порядок вищим,
        ніж передбачає $\Lambda$CDM на тому самому $z$.
  \item Загалом числова густина масивних галактик ($M_\star \gtrsim 10^{10}
        M_\odot$) при $z \gtrsim 8$ перевищує передбачення стандартних
        симуляцій (IllustrisTNG, EAGLE) в $10$--$100$ разів \cite{boylan_kolchin_2023}.
\end{itemize}

Суттєво, що ці аномалії проявляються саме на малих масштабах
$k \gtrsim 1\,h\,\mathrm{Mpc}^{-1}$~--- там де первинний спектр
$P(k)$ практично не обмежений Planck. Це і є вікно для нашої моделі.

%% --------------------------------------------------------
\subsection{Ідея}
%% --------------------------------------------------------

Якщо потенціал інфлатона має \textbf{сходинку}, м'яч сповільнюється~---
$\dot{\phi} \to 0$~--- і флуктуації що заморожуються в цей момент
мають більшу амплітуду згідно з~\eqref{eq:curvature_power_spec}.
Зламаний спектр:
\begin{equation}\label{eq:broken_spectrum}
  P(k) = A_s \left(\frac{k}{k_*}\right)^{n_s-1}
  \times \mathcal{F}(k;\, k_\mathrm{break},\, \Delta n),
\end{equation}
де функція переходу визначається як:
\begin{equation}\label{eq:F_def}
  \mathcal{F}(k;\, k_\mathrm{break},\, \Delta n)
  = \left[\frac{1}{2}\left(1 + \th\frac{\ln(k/k_\mathrm{break})}
    {\sigma}\right)\right]^{\!\Delta n},
\end{equation}
з параметром ширини переходу $\sigma$, який характеризує
«різкість» сходинки у потенціалі. У мінімальній версії моделі
$\sigma$ фіксується фізикою потенціалу і не є вільним
параметром MCMC; типові значення $\sigma \sim 0.5$--$1$
(у логарифмічній шкалі $k$) відповідають плавному переходу
протягом одного $e$-fold. Якщо форма потенціалу невідома,
$\sigma$ можна включити як третій параметр поряд з
$k_\mathrm{break}$ і $\Delta n$. Ця параметризація забезпечує:
\begin{itemize}
  \item $\mathcal{F} \to \left(\tfrac{1}{2}\right)^{\Delta n}$
        при $k \ll k_\mathrm{break}$~---
        стандартний спектр на великих масштабах, що бачить Planck
        (при $\Delta n \ll 1$ це $\approx 1$);
  \item $\mathcal{F} \to 1$ при $k \gg k_\mathrm{break}$~---
        підсилений спектр відносно великих масштабів;
  \item плавний перехід на масштабі $k_\mathrm{break}$~---
        без розривів у $P(k)$.
\end{itemize}
Модель має два нових параметри: $k_\mathrm{break}$ і $\Delta n$.
За $\Delta n > 0$ спектр на малих масштабах підсилюється відносно великих.

Більше потужності на малих масштабах $k \gtrsim 1\,\mathrm{Mpc}^{-1}$~---
гало тієї самої маси утворюються \textbf{раніше}~--- узгоджується з JWST.

%% --------------------------------------------------------
\subsection{Чому не суперечить CMB і $\sigma_8$}
%% --------------------------------------------------------

Planck чутливий до $k \lesssim 0.2\,\mathrm{Mpc}^{-1}$~--- злом не бачить.
$\sigma_8$ визначається масштабом $\sim 8\,\mathrm{Mpc}/h$, тобто
$k \sim 0.1$--$0.5\,\mathrm{Mpc}^{-1}$~--- майже не змінюється.

%% --------------------------------------------------------
\subsection{План роботи}
%% --------------------------------------------------------

\begin{enumerate}
  \item Огляд літератури: чи зроблено вже саме це поєднання.
  \item Мінімальна модель: 2 нових параметри $(k_\mathrm{break},\, \Delta n)$.
  \item CAMB/CLASS: розрахунок $P(k)$, $C_\ell$, $\sigma_8$.
  \item HMF (\texttt{hmf}, Python): порівняння з каталогами JWST.
  \item MCMC (Cobaya): posterior по повному набору параметрів.
  \item Стаття: чи є сумісна область параметрів для Planck + JWST + $\sigma_8$.
\end{enumerate}

%% --------------------------------------------------------
\subsection{Прогноз для майбутніх спостережень}
%% --------------------------------------------------------

Якщо MCMC знайде сумісну область параметрів
$(k_\mathrm{break}, \Delta n)$, модель дає конкретні
передбачення, перевірні наступним поколінням інструментів:

\begin{center}
\begin{tblr}{
  colspec = {llX},
  row{1}  = {font=\bfseries},
}
\toprule
Інструмент & Запуск & Що перевіряє \\
\midrule
Euclid \cite{euclid_2024}
  & 2023 (дані 2026+)
  & $P(k)$ галактик на $k \sim 0.1$--$5\,h/$Мпк,
    чутливий до $k_\mathrm{break}$ \\
Nancy Grace Roman (HLSS)
  & $\sim 2027$
  & weak lensing + galaxy counts при $z \sim 2$--$4$,
    обмежить $\Delta n$ \\
CMB-S4
  & $\sim 2029$
  & малі кутові масштаби CMB ($\ell \sim 3000$--$5000$),
    затухання Шовка \\
SKA (21~cm)
  & $\sim 2030$
  & $P(k)$ нейтрального водню при $z \sim 6$--$20$,
    прямий зонд $k_\mathrm{break}$ \\
\bottomrule
\end{tblr}
\end{center}

Ключовий тест: якщо $k_\mathrm{break} \sim 1$--$10\,h/$Мпк,
то Euclid і Roman матимуть достатню чутливість щоб
підтвердити або спростувати модель незалежно від CMB.
Комбінація «CMB + LSS + 21~cm» може обмежити
$k_\mathrm{break}$ з точністю $\sim 10\%$ і $\Delta n$
з точністю $\sim 0.05$ \cite{munoz_2022}.


\printbibliography

\end{document}